%% kapa: exam=1 round=28 year=2017 period=2 form=A subject=appraisal_law subject_ko=감정평가관계법규 seq=30 number=30
%% source: Q-Net 감정평가사 기출문제 (cst00309) artlSeq=5199603 file=extracted/1008/5199603/2017년 제28회 감정평가사 1차시험 문제지/2교시 A형.pdf
%% conversion: kordoc 4.24.51 -> tools/gen1.py (auto); review=proofread-sonnet; flags=none
\begin{question}{30}
공간정보의 구축 및 관리 등에 관한 법령상 토지의 이동 신청 및 지적정리 등에 관한 설명으로 옳은 것은?
\begin{choices}
  \item 토지소유자는 신규등록할 토지가 있으면 그 사유가 발생한 날부터 60일 이내에 지적소관청에 신규등록을 신청하여야 한다.
  \item 합병하려는 토지의 소유자가 서로 다른 경우에는 합병에 합의한 날로부터 90일 이내에 지적소관청에 합병을 신청하여야 한다.
  \item 지적소관청은 바다로 된 토지의 등록말소 신청을 하도록 통지받은 토지소유자가 그 통지를 받은 날로부터 60일 이내에 등록말소 신청을 하지 아니하면 등록을 말소한다.
  \item 지적소관청은 축척변경을 하려면 축척변경 시행지역의 토지소유자 과반수의 동의를 받아야 한다.
  \item 지적소관청은 합병하려는 토지가 축척이 다른 지적도에 각각 등록되어 있어 축척변경을 하는 경우에는 시ㆍ도지사의 승인을 받아야 한다.
\end{choices}
\end{question}
\answer{1}
